\documentclass[10pt]{amsart}
\usepackage[a4paper,margin=23mm]{geometry}
\usepackage{amsmath,amssymb,amsthm,mathtools,hyperref,graphicx,xspace,letltxmacro,xpatch}
\usepackage[english]{babel}
\usepackage{xfrac}

\newcommand{\ext}{\mathrm{ext}}

\usepackage{mathtools}
\usepackage{tikz}
\usetikzlibrary{arrows}

\newcommand\MTkillspecial[1]{% helper macro
\bgroup
\catcode`\&=9
\let\\\relax%
\scantokens{#1}%
\egroup
}
\newcommand\DeclarePairedDelimiterMultiline[3]{
\DeclarePairedDelimiter{#1}{#2}{#3}
\reDeclarePairedDelimiterInnerWrapper{#1}{star}{
\mathopen{##1\vphantom{\MTkillspecial{##2}}\kern-\nulldelimiterspace\right.}
##2
\mathclose{\left.\kern-\nulldelimiterspace\vphantom{\MTkillspecial{##2}}##3}}
}

\newcommand{\boH}{\mathcal{H}}
\newcommand{\boC}{\mathcal{C}}
\newcommand{\boI}{\mathcal{I}}
\newcommand{\boB}{\mathcal{B}}
\newcommand{\boD}{\mathcal{D}}
\newcommand{\boT}{\mathcal{T}}
\newcommand{\Z}{\mathbb{Z}}
\newcommand{\N}{\mathbb{N}}
\newcommand{\R}{\mathbb{R}}
\newcommand{\CC}{\mathbb{C}}
\newcommand{\Sbb}{\mathbb{S}}
\newcommand{\dd}{\mathop{}\!\mathrm{d}}
\newcommand{\ic}{\mathbf{i}}
\newcommand{\bpi}{\boldsymbol{\pi}}
\DeclarePairedDelimiterMultiline{\abs}{\lvert}{\rvert}
\DeclarePairedDelimiterMultiline{\norm}{\lVert}{\rVert}
\DeclarePairedDelimiterMultiline{\pare}{(}{)}
\newcommand{\pareB}[1]{\Bigl(#1\Bigr)}
\newcommand{\pareg}[1]{\bigl(#1\bigr)}
\newcommand{\paregg}[1]{\biggl(#1\biggr)}
\newcommand{\normB}[1]{\Bigl\lVert#1\Bigr\rVert}
\newcommand{\normg}[1]{\bigl \lVert#1\bigr\rVert}
\newcommand{\normgg}[1]{\biggl \lVert#1\biggr\rVert}
\newcommand{\absB}[1]{\Bigl\lvert#1\Bigr\rvert}
\newcommand{\absg}[1]{\bigl \lvert#1\bigr\rvert}
\newcommand{\absgg}[1]{\biggl \lvert#1\biggr\rvert}
\newcommand{\st}{\::\:}
\newcommand{\restr}{|}

\DeclareMathOperator{\ind}{ind}
\DeclareMathOperator{\Fix}{Fix}
%\DeclareMathOperator{\deg}{deg}
\DeclareMathOperator{\sgn}{sgn}
\DeclareMathOperator{\Id}{Id}
\DeclareMathOperator{\Card}{Card}
\DeclareMathOperator{\Leb}{Leb}
\DeclareMathOperator{\dLeb}{dLeb}
\DeclareMathOperator{\Span}{Span}
\newcommand{\coloneqq}{\mathrel{\mathop:}=}
\DeclareMathOperator{\flattr}{\mathrm{tr}^{\flat}}
\DeclareMathOperator{\tr}{tr}

\let\epsilon\varepsilon
\let\oldphi\phi
\let\phi\varphi

\newtheorem{thm}{Theorem}[section]
\newtheorem{prop}[thm]{Proposition}
\newtheorem{definition}[thm]{Definition}
\newtheorem{lem}[thm]{Lemma}
\newtheorem{cor}[thm]{Corollary}

\theoremstyle{definition}
\newtheorem{ex}{Example}
\newtheorem{rmk}[thm]{Remark}

\numberwithin{equation}{section}


\usepackage{xfrac,xcolor,graphicx}\graphicspath{{./figures/}}\providecommand{\skpt}{\par\smallskip\noindent}\providecommand{\pointir}{.\ }\providecommand{\eg}{e.g.}\providecommand{\ie}{i.e.}\providecommand{\resp}{resp.}\providecommand{\cf}{cf.}\providecommand{\etc}{etc.}\let\Subsection\subsection\let\Subsubsection\subsubsection
